\section{Finite recognition and protected geometric tests}\label{sec:recognition}
We distinguish a statistical reconstruction of a physical body from its
translation class in the unknown periodic table. Spatial clustering in
Lemma~\ref{lem:cloud} solves the first problem locally. The following
finite descriptor solves the second without extending a contact germ.

\subsection{Support fingerprints in the normal frame}
For an oriented shortest bridge, translate its source contact to zero
and rotate the normal towards the target to $(1,0)$. The transverse axis
has either of two signs. Write $p_0^e,p_1^e$ for the support functions of
the two complete bodies in this one frame, without separate recentering.
Their absolute values and first derivatives are bounded by $K_0$ in
\eqref{eq:fingerprint-size}. Define
\begin{equation}\label{eq:descriptor}
 F(e)=\bigl(g_e,(p_0^e(2\pi j/m))_{j=0}^{m-1},
                    (p_1^e(2\pi j/m))_{j=0}^{m-1}\bigr).
\end{equation}
Reflection across the normal permutes $j$ to $-j$ modulo $m$ in both
blocks. It is a coherent reversal of the pair, not separate endpoint signs.

\begin{lemma}[Locking physical copies from complete pairs]\label{lem:locking}
Suppose two normal-frame pairs from one table have support differences
at most $\Delta_0$ on the full circle. Then they represent the same
oriented bridge up to translation.
\end{lemma}
\begin{proof}
If two support functions differ by at most $\Delta$, their Steiner
points differ by at most $2\Delta$, since
$s(C)=\pi^{-1}\int_0^{2\pi}p_C(\theta)n(\theta)\dd\theta$.
Their centered support functions consequently differ by at most
$3\Delta$ and by at most $9\Delta$ in $L^2$. The shape gap excludes
different source types when $\Delta\le\Delta_0$. The exact congruence
between the two source copies, expressed in their normal frames, has
orthogonal part $O$. The reflection margin excludes its orientation-reversing
component; the rotation margin bounds its angle by $9\Delta/\eta$.

Align the two physical source copies by their true lattice translation.
The preceding centering and angle estimates show that their target
images then differ in Hausdorff distance by at most
\[
 \{3+9(R_++2D_0)/\eta\}\Delta=C_0\Delta<d_0.
\]
For completeness, the centering displacement costs $2\Delta$; rotation
about a source Steiner point costs at most
$(R_++2D_0)9\Delta/\eta$ on the pair, and the original target comparison
costs $\Delta$. Distinct target bodies in that physical table have
separation at least $d_0$, so the two targets coincide. The closest
segment of two disjoint strictly convex bodies is unique. It therefore
fixes both contacts, the normal and, by the reflection exclusion, the
transverse sign. This proves the assertion.
\end{proof}

\begin{proof}[Proof of Theorem~\ref{thm:fingerprint-main}]
If the node differences are at most $\chi=\Delta_0/4$, the derivative
bound gives, at any circle point,
\[
 |p_b^e-p_b^{e'}|\le\chi+2K_0\pi/m
                          \le\Delta_0/4+\Delta_0/8<\Delta_0.
\]
Here $\pi<4$ and $m\ge64K_0/\Delta_0$ were used. Apply
Lemma~\ref{lem:locking}. The contrapositive gives separation greater
than $\chi$ for inequivalent oriented classes. All ingredients of
\eqref{eq:fingerprint-size} are finite arithmetic in numerical margins.
With each estimated descriptor within $\chi/16$ of its truth, same-class
distances are at most $\chi/8$, and different-class distances exceed
$7\chi/8$. A threshold $\chi/2$, testing both coherent reversals,
therefore recognizes the class and sign. Fixed-order support, normal
and gap errors from Theorem~\ref{thm:probe-main} make these comparisons
valid once $C\nu<\chi/16$.
\end{proof}

This polynomial-size fingerprint uses whole-body support values produced
by bounded-range sampling. It is not a shorter list of values on the
same short contact interval used in the preceding analytic theorem.
That distinction explains why its separation needs no continuation
constant. It also prevents a false claim that a smooth body is globally
determined by a finite contact jet.

\subsection{One explicit normal-root certificate}\label{sec:root}
We give quantitative choices for the local root and clearance tests.
In true unit-speed charts near a closest pair, $D(s,u)=|x(s)-y(u)|$
satisfies
\begin{equation}\label{eq:normal-hessian}
 \nabla D=0,\qquad D''=
 \begin{pmatrix}g^{-1}+\kappa_0&-g^{-1}\\
                -g^{-1}&g^{-1}+\kappa_1\end{pmatrix}.
\end{equation}
Its least eigenvalue is at least $\mu=\kappa_-$. In unit-speed charts
one may use the explicit Hessian-Lipschitz bound
\[
 L_D=64\{1+d_0^{-2}+\kappa_+/d_0+\kappa_+^2
                                      +2M_6\kappa_+^3\}.
\]
Indeed, writing $\rho=p+p''$, one has
$|\dd\kappa/\dd s|=|\rho'|/\rho^3\le2M_6\kappa_+^3$.
The first three derivatives of a unit-speed arc are bounded by
$1$, $\kappa_+$ and $2M_6\kappa_+^3+\kappa_+^2$.
The first three derivatives of Euclidean distance at separation at least
$d_0$ are bounded by $1,d_0^{-1},6d_0^{-2}$. The multivariate chain
rule gives the conservative displayed bound.
Set
\begin{equation}\label{eq:root-radius}
 r=\min\{r_c/8,d_0/32,\mu/(32L_D)\}.
\end{equation}
The initial encountered patches are accepted only inside a protected
neighborhood on which the true Hessian has least eigenvalue at least
$3\mu/4$. This condition is verifiable with the cloud's derivative
enclosures; genuine witness returns have uniform slack.

Here is a specific certificate. At a trial center $z$, let $B$ be a
symmetric approximation to $D''(z)$ with error at most $\mu/16$ and
least eigenvalue at least $\mu/2$. Let $\epsilon_g$ bound the gradient-approximation error uniformly
on the ball, as provided by the cloud enclosures. Verify
the inequality
\begin{equation}\label{eq:newton-test}
 \|\widehat{\nabla D}(z)\|+\epsilon_g\le\mu r/8,
 \qquad \|D''(x)-B\|\le\mu/4\quad(x\in\overline B(z,r)).
\end{equation}
The second condition follows from $L_Dr+\mu/16\le\mu/4$ when the
ball stays in the certified rectangle. The map
$x\mapsto x-B^{-1}\nabla D(x)$ has Lipschitz constant at most $1/2$
on this ball and moves its center by at most $r/4$. It maps the ball
strictly into itself and has a unique fixed point. The root lies within
$2\|B^{-1}\|\epsilon_g$ of the corresponding root for the central
approximation, with the approximation residual added to $\epsilon_g$
when it is not zero. This gives finite root enclosures, rather than an
unstated interval-Newton call.

Approximate roots can be found by a finite mesh followed by contraction,
or by minimizing distance between the reconstructed strictly convex
bodies. The uniform Hessian and chart slack make the tests pass for all
sufficiently accurate clouds of a protected bridge. Refinement is capped
at the deterministic sample and arithmetic budgets; indeterminate inputs
are rejected. Convex supporting half-planes show that a certified facing
normal is the global shortest segment for those two bodies.

\subsection{Clearance and record assignment}
Let $C_0^e,C_1^e$ denote the two \emph{physical bodies} defining the
candidate bridge; different copies of the same type are still different
physical bodies. Define
\[
 \mathcal O_e^{\rm other}
  =\bigcup_{C\text{ a physical body},\ C\ne C_0^e,C_1^e} C.
\]
Delete exactly the two corresponding spatial cloud components from
$P_{\rm all}$, not every component with either endpoint's shape. The
argument for \eqref{eq:distance-cloud}, applied after this deletion,
gives the same capped-distance enclosure for $\mathcal O_e^{\rm other}$.
The padded cloud region must contain every body meeting the capped
neighborhood of the entire tested segment. Partial outer components are
kept, just as in Section~\ref{sec:distance}.

Put $\bar c=\min(c_0,d_0/4)$ and use cap at least $\bar c$. On the
closed candidate segment, a mesh of spacing at most $\bar c/8$, total
distance and root-position enclosure radius at most $\bar c/8$, and
reported distances at least $3\bar c/4$ imply
\[
 \dist([x_0,x_1],\mathcal O_e^{\rm other})\ge\bar c/2.
\]
This follows directly from the one-Lipschitz property of distance.
Protected bridges with the stipulated clearance pass at sufficiently
small error. The endpoint bodies have deliberately not been included
in this inequality.

Their contribution is separate and exact. Write $x_1=x_0+gn$ for the
certified common normal. The facing supporting half-planes give, for
$0\le s\le g$,
\[
 \dist(x_0+sn,C_0^e)=s,\qquad
 \dist(x_0+sn,C_1^e)=g-s.
\]
Choose $a=\min(d_0/8,g/4)$ and treat $s\in[0,a]$ and
$s\in[g-a,g]$ in the protected endpoint charts, restricted to the
exterior sides of the supports. On $[a,g-a]$ a tube of radius smaller
than $\min(a/2,\bar c/4)$ avoids all solids. Root uncertainty is included
in the enclosure budget above. An upper bound on $a$ is not used as a
lower bound on endpoint distance; the explicit tube radius uses $a$
itself. All distances are computed from this same cloud, not obtained
from a separate geometric sensor.

At most one record is extracted from a trace, and it must have exactly
three impacts before its recorded horizon. First and third impacts must belong to the same spatial cloud component,
with the middle impact on the other body. Inside fixed protected patches,
the stationary two-flight branch is unique. Record assignment additionally
checks the collision order, gap cutoff below $R_+$, protected patches,
twist and adjacent-flight margins. Every trajectory in a common smaller
return rectangle of a clear bridge of length below $R_-$ passes these
checks on the confidence event. The slack $R_+-R_-$ protects the cutoff.
For a registered branch, future trajectories are classified with the
same support descriptor and assigned intrinsic coordinates using the
reconstructed charts and their true-error enclosures. Interior histogram
rectangles have a fixed collar. Ambiguity in a bin is consequently only
an endpoint-coordinate band of thickness $C\nu$, not an uncharged change
of branch. Failed or out-of-rectangle trajectories remain outcomes.

This proves the record part of Theorem~\ref{thm:probe-main}. Its confidence
event is simultaneous for all geometry extracted from one cloud. It does
not say that each statistically accepted record is infallible. The next
two sections explain the finite-list certificate and how independent
revalidation makes the distinction harmless for an anytime guarantee.
